\vssub
\subsection{~运行测试案例} \label{sec:tests}
\vssub

如果 \ws{} 已成功安装并编译，可以通过交互式运行不同程序元素来测试它。为此，创建一个名为 {\file work} 的顶层目录，并复制 {\file model/inp} 或 {\file model/nml} 目录中的输入文件。现在不再提供通用 switch 文件，但可以从 {\file bin} 目录中选择一个 switch 文件，然后测试示例输入文件。因此，应当可以通过输入以下命令运行所有模式元素：
\command{ww3\_grid | more \\
  ww3\_strt | more \\
  ww3\_bound | more \\
  ww3\_prep | more \\
  ww3\_shel | more \\
  ww3\_outf | more \\
  ww3\_outp | more \\
  ww3\_ounf | more \\
  ww3\_ounp | more \\
  ww3\_trck | more \\
  ww3\_grib | more \\
  gx\_outf | more \\
  gx\_outp | more }
其中添加 {\code more} 命令是为了允许在屏幕上检查输出。这个 {\code | more} 可以替换为重定向到输出文件，例如 \command{ww3\_grid > ww3\_grid.out }。注意，只有链接了用户提供的打包例程时，{\code ww3\_grib} 才会提供 GRIB 输出。还需注意，不提供用于 {\file ww3\_multi} 的简单交互式测试案例。随后可从 work 目录运行 GrADS，为这些计算生成图形输出。所有中间输出文件都放在 {\file work} 目录中，可通过输入 \command{w3\_clean} 方便地删除。

直到 3.14 版，\ws{} 都随附一组简单测试，用于确定并评估模式基本功能是否行为正确。在下一版模式的早期开发中，Erick Rogers 和 Tim Campbell 将这些测试转换为更易以自动化方式运行的回归测试。直到模式 4.06 版，这些修改后的测试收集在 {\file nrltest} 目录中，同时旧测试保留在 {\file test} 目录中。在模式 4.07 版，{\file nrltest} 被采纳为 \ws{} 的新测试案例，并放入新的 {\file regtests} 目录；最终，{\file test} 中剩余的真实世界测试案例被移到 {\file cases} 目录，同时完全停用 {\file test} 目录。{\file regtests} 目录中提供以下回归测试。

\begin{flist}
\fit{ww3\_tp1.1}{沿赤道绕全球的一维传播（无陆地）。}
\fit{ww3\_tp1.2}{沿经线的一维传播（无陆地）。}
\fit{ww3\_tp1.3}{一维传播，浅化测试。}
\fit{ww3\_tp1.4}{一维传播，谱折射（$x$）。}
\fit{ww3\_tp1.5}{一维传播，谱折射（$y$）。}
\fit{ww3\_tp1.6}{一维传播，海流导致的波浪阻挡。}
\fit{ww3\_tp1.7}{一维传播，IG 波生成。}
\fit{ww3\_tp1.8}{一维传播，海滩上的波浪破碎。}
\fit{ww3\_tp1.9}{一维传播，Beji and Battjes (1993) 沙坝水槽案例。}
\fit{ww3\_tp2.1}{相对于网格成角度的二维传播。}
\fit{ww3\_tp2.2}{无陆地半球上的二维传播（带方向扩展）。}
\fit{ww3\_tp2.3}{二维传播，GSE 测试。}
\fit{ww3\_tp2.4}{二维传播，东太平洋曲线网格测试。}
\fit{ww3\_tp2.5}{二维传播，北极网格、曲线网格测试。}
\fit{ww3\_tp2.6}{二维传播，Limon Harbor 非结构网格测试。}
\fit{ww3\_tp2.7}{二维非结构网格上的反射。}
\fit{ww3\_tp2.8}{二维规则网格上的潮汐分潮。}
\fit{ww3\_tp2.9}{障碍物网格测试。}
\fit{ww3\_tp2.10}{SMC 网格测试。}
\fit{ww3\_tp2.11}{旋转网格测试。}
\fit{ww3\_tp2.12}{波系跟踪测试。}
\fit{ww3\_tp2.13}{相对于网格成角度的传播测试（三极网格）。}
\fit{ww3\_tp2.14}{使用 OASIS 耦合器的玩具模式测试。}
\fit{ww3\_tp2.15}{时空极值参数测试。}
\fit{ww3\_tp2.16}{采用北极极点处理的 SMC 网格二维传播。}
\fit{ww3\_tp2.17}{非结构网格，Card Deck 与显式/隐式格式区域分解对比；使用 ww3\_gint 进行网格整合；grib2 输出（ww3\_grib）。}
\fit{ww3\_tp2.18}{ww3\_prnc/ww3\_prtide 中潮汐分潮测试。}
\fit{ww3\_tp2.19}{带区域分解/隐式格式的非结构网格，深度破碎和三波源项采用 Neumann 边界条件（Boers 1996）。}
\fit{ww3\_tp2.20}{植被源项（Anderson and Smith 2012）。}
\fit{ww3\_tp2.21}{全球非结构网格（周期性和障碍物）；使用 ww3\_gint 进行网格整合；ww3\_grib 输出。}
\fit{ww3\_ts1  }{源项测试，受时间限制的增长。}
\fit{ww3\_ts2  }{源项测试，受风区限制的增长。}
\fit{ww3\_ts3  }{源项测试，带单个移动网格的飓风。}
\fit{ww3\_ts4  }{源项测试，未解析障碍物。}
\fit{ww3\_tic1.1}{波-冰相互作用，$S_{ice}$ 的一维测试。}
\fit{ww3\_tic1.2}{波-冰相互作用，“浅化”效应的一维测试。}
\fit{ww3\_tic1.3}{波-冰相互作用，折射效应的一维测试。}
\fit{ww3\_tic1.4}{波-冰相互作用，带浮冰和冰厚的一维测试。}
\fit{ww3\_tic2.1}{波-冰相互作用，$S_{ice}$ 的二维测试。}
\fit{ww3\_tic2.2}{波-冰相互作用，带非均匀冰的二维测试。}
\fit{ww3\_tic2.3}{波-冰相互作用，带厚度递增的均匀冰的二维测试。}
\fit{ww3\_tbt1.1}{波-泥相互作用，$S_{mud}$ 的一维测试。}
\fit{ww3\_tbt2.1}{波-泥相互作用，$S_{mud}$ 的二维测试。}
\fit{ww3\_tpt1.1}{替代谱分区方法测试。}
\fit{ww3\_ta1   }{ww3\_uprstrt，更新均匀条件的重启动文件（1 点模式）。}
\fit{mww3\_test\_01}{湿干等条件下扩展网格掩膜测试。}
\fit{mww3\_test\_02}{带单个内嵌网格的双向嵌套测试。}
\fit{mww3\_test\_03}{重叠网格和双向嵌套测试（高分辨率网格中带海滩的 6 网格版本）。}
\fit{mww3\_test\_04}{静止涌浪条件下带双向嵌套的海流或海山测试。}
\fit{mww3\_test\_05}{带移动网格的三个嵌套飓风网格测试。}
\fit{mww3\_test\_06}{使用 {\file ww3\_multi} 的不规则网格测试。}
\fit{mww3\_test\_07}{使用 {\file ww3\_multi} 的非结构网格测试。}
\fit{mww3\_test\_08}{带风和冰输入的测试。}
\fit{mww3\_test\_09}{等等级 SMC 子网格测试。}
\fit{ww3\_ufs1.1}{全球 1 度网格测试，使用 ww3\_multi/ww3\_multi\_esmf；风/流/冰输入；使用 ww3\_gint 进行网格整合；grib2 输出（ww3\_grib）；重启动/MPI 任务/线程可复现性（ufs app）。}
\fit{ww3\_ufs1.2}{GFSv16 设置测试，使用 ww3\_multi/ww3\_multi\_esmf；风/流/冰输入；使用 ww3\_gint 进行网格整合；grib2 输出（ww3\_grib）；重启动可复现性（ufs app）。}
\end{flist}

\noindent
这些回归测试现在使用 {\file regtests/bin} 目录中的 {\file run\_test} 脚本运行（主要作者：Tim Campbell）。如何运行该脚本，包括选项，可通过运行 \command{run\_test -h} 查看。运行该命令的输出如图~\ref{fig:runtest} 所示。测试案例存储在 {\file regtests} 目录下的各个目录中，例如 {\file regtests/ww3\_tp1.1}。例如，{\file /ww3\_tp1.1} 的内容可能为：

\begin{flist}
\fit{info }{包含测试案例信息的文件。}
\fit{input}{包含测试案例输入文件的永久目录。}
\fit{work\_PR3}{模式输出的暂存目录（在本例中，文件名如此是因为用户指定了 ``{\file run\_test -w work\_PR3 ...''}）。}
\end{flist}

\noindent
现在还提供了一个回归测试矩阵，代码开发者使用它来确保新模式版本不会破坏旧模式版本。该矩阵的核心是文件 {\file regtests/bin/matrix.base}。{\file regtests/bin/matrix\_zeus\_HLT} 给出了运行方式示例，这是 Hendrik 在 NCEP Zeus R\&D 计算机上运行该矩阵的驱动脚本\footnote{~请以此为蓝本为自己的设置构建自己的驱动脚本，而不是编辑该文件。}。若要运行它，在 {\file regtests} 目录中创建指向它的链接，并在脚本中设置所需选项标志后执行。这样会在 {\file retests} 中生成一个 {\file matrix} 文件，随后可按需要交互式或批处理运行。若有需要，也可以进一步手工编辑该文件。{\file regtests} 下的 {\file bin} 目录包含以下工具。

\begin{flist}
\fit{cleanup          }{清理工作目录。}
\fit{comp\_switch     }{比较测试案例内部和测试案例之间的开关。{\file comp\_switch -h} 提供文档。}
\fit{matrix.base      }{生成测试案例矩阵的核心脚本。}
\fit{matrix.comp      }{比较分别检出的模式版本之间测试案例矩阵输出的脚本。}
\fit{matrix\_zeus\_HLT}{用于 {\file matrix.base} 的驱动脚本示例。}
\fit{run\_test        }{如上所述的基本测试脚本。}
\end{flist}

\noindent
注意，高效运行回归测试矩阵需要尽量减少各回归测试之间重新编译代码的需求。这通过 {\file matrix.base} 中回归测试的排序实现。确保相同 switch 文件被识别为相同文件的一种方法，是对它们进行系统排序。可用主 {\file bin} 目录中的脚本 {\file sort\_switch} 完成。该脚本会添加缺失开关的默认值，也可用于从文件中删除或添加开关。运行 \command{comp\_switch -h} 查看该脚本文档。

\vspace{\baselineskip} \noindent
最后，{\file cases} 目录保存如下真实世界测试案例。

\begin{flist}
\fit{mww3\_case\_01}{以 Trondheim 为重点的五网格大西洋案例。}
\fit{mww3\_case\_02}{以 Alaska 为重点的三网格太平洋案例。}
\fit{mww3\_case\_03}{NCEP 曾用作全球模式的原始多网格案例。}
\end{flist}

\noindent
这些案例中的每一个都是执行完整模式运行的单个脚本。执行脚本之前，请使用脚本开头文档中指明的开关编译模式。这些脚本使用的附加数据包含在以下目录中：

\begin{flist}
\fit{mww3\_data\_00}{所有示例案例使用的风场和冰数据。}
\fit{mww3\_data\_{\it nn}}{脚本 {\file mww3\_case\_{\it nn}} 所需的特定数据。}
\end{flist}

\noindent
这些示例可作为建立其他真实模式应用的蓝本。

\begin{figure}
{\scriptsize
\IfFileExists{run_test.out}
  {\input{run_test.out}}
  {\noindent\texttt{run\_test.out} 未生成。运行 \texttt{make inp} 可生成此帮助输出。}
}
\caption{使用 {\file -h} 命令行选项运行 {\file run\_test} 得到的选项。} \label{fig:runtest}
%\botline
\end{figure}

